\documentclass[11pt]{article}
\usepackage[review]{acl}
\usepackage{times}
\usepackage{latexsym}
\usepackage[T1]{fontenc}
\usepackage[utf8]{inputenc}
\usepackage{microtype}
\usepackage{booktabs}
\title{Adapting Pretrained Language Models to State--Prediction Separation}
\author{Anonymous Authors}
\begin{document}
\maketitle
\begin{abstract}
Can an existing autoregressive model be converted to a state--prediction separated (SPS) Transformer without losing quality? Two-stage adaptation recovers much of the initial conversion loss across five base models, but no final CE or CE+KD checkpoint preserves every measured source metric. For an instruction-tuned model, improved accuracy on a shared partial MMLU subset coexists with worse instruction following. A separate single-request serving pilot finds shared-backbone SPS speculative decode-throughput gains of 1.58--1.80$\times$ with target decode CUDA graphs, without establishing preserved sampled quality. These preliminary results motivate treating language modeling, instruction retention and serving efficiency as separate adaptation objectives.
\end{abstract}

\section{Question and approach}
SPS separates persistent text states from next-token prediction states \citep{monea2026sps}. We study conversion of pretrained weights rather than training an SPS model from scratch. In zero-window SPS, a prediction slot follows each text token; text positions cannot attend to prediction slots. Each prediction slot reads the causal text prefix and itself.

We first learn only the new prediction embedding, then unfreeze the model. We compare cross-entropy (CE) with CE plus distribution matching to the frozen source model (CE+KD). Base experiments cover SmolLM2-135M, 360M and 1.7B, and Qwen3-0.6B-Base and 1.7B-Base. Instruction experiments use post-trained Qwen3-1.7B and matched AR continuation controls. Implementation and detailed results are deferred to the appendix.

\section{Main observations}
\paragraph{Adaptation repairs conversion, but retention is incomplete.}
For SmolLM2-135M, untrained conversion raises held-out NLL from 2.828 to 6.165; two-stage CE reduces it to 2.857. The staged recipe improves LAMBADA over immediate full adaptation by 0.80 percentage points, but remains below the source. Across all five base models, KD improves final NLL over CE, yet neither recipe preserves every source benchmark and NLL (Appendix~\ref{app:base}).

\paragraph{Knowledge-test gains do not imply instruction retention.}
Table~\ref{tab:main} compares the instruction-tuned source and final constant-LR adaptations. SPS CE improves partial generative MMLU accuracy while losing 12.93 points on non-thinking IFEval \citep{zhou2023ifeval}. AR continuation also loses non-thinking IFEval quality, so the decline is not unique to SPS. However, SPS remains below its matched AR control in that mode.

\begin{table}[t]
\centering\small
\begin{tabular}{lrrr}
\toprule
& Partial & \multicolumn{2}{c}{IFEval}\\
Model & MMLU T & T & NT\\
\midrule
Source & 67.22 & 65.62 & 65.43\\
SPS CE & \textbf{70.74} & 64.14 & 52.50\\
SPS CE+KD & 69.06 & 64.33 & 55.82\\
AR CE & 69.94 & \textbf{66.91} & 54.16\\
AR CE+KD & 67.95 & 60.26 & 60.81\\
\bottomrule
\end{tabular}
\caption{Qwen3-1.7B instruct, step 512 adaptations. Percent accuracy; IFEval uses prompt-level strict accuracy on all 541 prompts. T/NT: thinking/non-thinking. MMLU uses the same 3,859 questions for every row, not the full test.}
\label{tab:main}
\end{table}

The shared MMLU subset \citep{hendrycks2021mmlu} covers 27.48\% of the test and is an ordered subject prefix, not a random sample. Invalid answers count as errors. Their rate falls from 19.93\% to 11.27\% for SPS CE, so the accuracy gain need not reflect knowledge alone. We do not claim full-MMLU retention.

\paragraph{More updates are not sufficient.}
Neither warmup/cosine scheduling nor continuation on the same instruction corpus consistently removes the IFEval deficit. Thinking targets are source-generated, whereas non-thinking targets come from SmolTalk; this confounds generation mode with target provenance. Even CE therefore includes sequence-level distillation for thinking. Matching target provenance is a more direct next test than simply extending training.

\paragraph{Serving gains depend on the verifier and execution mode.}
In a separate vLLM pilot on 100 MMLU STEM questions, a jointly adapted SPS model uses one backbone for multi-token proposals and verification. With target decode CUDA graphs, SPS-MTP reaches 1.58$\times$ thinking and 1.80$\times$ non-thinking decode-throughput gains against the same final SPS weights without speculation. EAGLE-3 reaches 1.46$\times$ and 1.60$\times$ against the original AR verifier. These are separate baselines, not a fixed-target comparison between draft methods. SPS-MTP thinking accuracy is 63\%, versus 70\% for its SPS baseline; graph execution itself leaves all sampled token sequences unchanged from eager execution. Appendix~\ref{app:serving} gives the protocol and both execution-mode results. This STEM pilot is distinct from the larger partial MMLU comparison above.

\section{Takeaway}
Pretrained SPS conversion is feasible as an adaptation procedure, but lossless quality transfer is not established. Distribution matching helps base-model NLL; instruction following remains a separate failure mode. Shared-backbone speculation offers measured serving gains on the small pilot, while quality retention remains unresolved.

\section*{Limitations}
This is an exploratory, single-training-seed draft. Partial MMLU is not representative of the full test; instruction-score differences have no uncertainty estimates here. Most base configurations lack AR continuation controls. The speed results cover only a small, single-GPU, batch-one pilot with target decode graphs and no Torch compilation; they do not establish production throughput, fixed-length speedups or speculative sampling-distribution equivalence. No capacity claim is made. Repeated training seeds and broader matched quality and serving evaluations are needed before submission.

\bibliography{references}
\clearpage
\appendix
\section{Base-model retention}\label{app:base}
We define observed retention as no decrease in any selected benchmark and no increase in held-out NLL. This descriptive rule is not a statistical equivalence test. Table~\ref{tab:nll} reports final NLL; compare recipes within a model, not token-level losses across tokenizers. All final CE and KD variants fail the full retention rule.

\begin{table}[h]
\centering\small
\begin{tabular}{lrrr}
\toprule
Model & Source & CE & CE+KD\\
\midrule
SmolLM2-135M & 2.828018 & 2.856920 & 2.851728\\
SmolLM2-360M & 2.605900 & 2.617682 & 2.614695\\
SmolLM2-1.7B & 2.305988 & 2.326675 & 2.321363\\
Qwen3-0.6B & 2.754853 & 2.781615 & 2.761109\\
Qwen3-1.7B & 2.513415 & 2.531043 & 2.519932\\
\bottomrule
\end{tabular}
\caption{Base-model held-out NLL (nat/token); lower is better.}
\label{tab:nll}
\end{table}

SmolLM2-135M embedding-only adaptation reaches LAMBADA 42.81\% and NLL 3.049, versus source 42.71\% and 2.828. Recovering one metric does not establish overall retention. Two-stage CE reaches LAMBADA 41.90\%; its paired difference from immediate full tuning is +0.80 points [0.29, 1.30], while its difference from the source is $-0.82$ [$-1.59$, $-0.02$]. Intervals are per-metric 95\% intervals over evaluation examples, not training seeds or simultaneous guarantees.

Qwen3-1.7B CE raises ARC-Easy from 68.52\% to 73.57\%, but lowers HellaSwag from 66.37\% to 64.67\%, SciQ from 94.90\% to 93.60\%, and LAMBADA from 63.09\% to 61.69\%. KD partially recovers these losses to 65.59\%, 94.50\%, and 62.25\%, respectively. Base AR continuation and immediate-unfreezing ablations were run for the 135M pilot, not the full model matrix.

\section{Instruction data and ablations}
The instruction model is Qwen3-1.7B, not Qwen3-1.7B-Base. Adaptation uses assistant-only loss on SmolTalk-derived dialogs. Thinking responses are generated by the frozen source; non-thinking responses are corpus targets. Thus CE versus CE+KD isolates additional distribution matching, not the presence versus absence of teacher-generated text. Neither the AR controls nor the mode gap alone establishes a causal explanation for the deficit.

Table~\ref{tab:ablations} reports full IFEval for the SPS arms. Continuation reuses the same corpus; it tests additional updates, not more unique data. Intermediate checkpoints and schedule variants are exploratory comparisons rather than a separately validated checkpoint-selection rule.

\begin{table}[h]
\centering\small
\begin{tabular}{llrr}
\toprule
Recipe & Endpoint & T & NT\\
\midrule
CE & Constant, 512 & 64.14 & 52.50\\
CE+KD & Constant, 512 & 64.33 & 55.82\\
CE & Cosine, 512 & 64.51 & 48.98\\
CE+KD & Cosine, 512 & 62.48 & 59.33\\
CE & Continued, 2048 & 60.81 & 53.79\\
CE+KD & Continued, 2048 & 57.86 & 52.87\\
\bottomrule
\end{tabular}
\caption{SPS instruction-following ablations: full IFEval prompt-level strict accuracy (\%).}
\label{tab:ablations}
\end{table}

\section{Partial MMLU and missing evaluations}
Thinking runs did not all finish. We compare the intersection across the source and twelve adaptations: CE/KD for SPS/AR at constant-LR step 384, constant-LR step 512 and cosine step 512. The 3,859 shared questions cover 25 of 57 subjects, with the last subject incomplete. Question hashes, targets, subjects and generation seeds agree across variants; manifests agree on generation settings and dataset revision. No unfinished JSONL rows were encountered. Truncated generations remain in the denominator and are scored by the extracted final answer; unparseable answers are errors.

\begin{table}[h]
\centering\small
\begin{tabular}{llrr}
\toprule
Schedule/step & Model & Acc. & Invalid\\
\midrule
Source & AR & 67.22 & 19.93\\
Constant/384 & SPS CE & 69.45 & 14.23\\
 & SPS KD & 68.26 & 15.73\\
 & AR CE & 67.27 & 17.34\\
 & AR KD & 68.36 & 16.53\\
Constant/512 & SPS CE & 70.74 & 11.27\\
 & SPS KD & 69.06 & 14.43\\
 & AR CE & 69.94 & 12.80\\
 & AR KD & 67.95 & 16.82\\
Cosine/512 & SPS CE & 67.92 & 16.30\\
 & SPS KD & 69.19 & 16.20\\
 & AR CE & 69.71 & 14.67\\
 & AR KD & 68.54 & 17.05\\
\bottomrule
\end{tabular}
\caption{Thinking MMLU on the identical 3,859-question subset (\%). KD denotes CE+KD.}
\end{table}

There are no saved non-thinking MMLU responses for the source or three constant-step-512 adaptations. The all-variant non-thinking intersection is therefore empty. Nine available variants share 1,377 questions across 11 subjects, but this cannot establish retention against the source. Absolute thinking and non-thinking scores on these different subsets are not comparable. The shared-subset result is descriptive and does not establish statistical significance.

\section{Evidence and draft scope}
The code registries are the single source of experimental configuration. Base runs use \texttt{scripts/ar\_sps\_config.py}; instruction, scheduling and continuation runs use the corresponding \texttt{instruct}, \texttt{sched} and \texttt{extend} configuration modules. Training and evaluation workers are separate. Measured evidence locations are documented in the accompanying draft README and Russian appendix notes.

This draft excludes MTP, serving speed and the long base-training follow-up. The remaining full benchmark tables, complete model/data bibliography and publication-ready uncertainty analysis are not yet incorporated. No new experiments were run for manuscript preparation.

\end{document}
